% Transcription — fonds Alexandre Grothendieck, shelfmark 50, batch 1 (pages 1-10).
% Established from the facsimile archives/batches/50/batch-01.pdf (a mirror of
% https://grothendieck.umontpellier.fr/50.pdf), per the skill
% transcribe-grothendieck. The folder is ten pages and this is its only batch.
%
% Page 1 is a grey wrapper carrying, in pencil in the top right corner, the four
% words "Sur bouquin Lang" and nothing else: it takes no \page{} and its words
% open the file as the section heading. Pages 3 and 5 are the recto and verso of
% a Bourbaki typescript on valuations, in another hand and unrelated to the
% notes; they are simply absent, and the gap in the numbering is the record.
%
% The eight remaining pages are two distinct things.
% - Pages 2 and 4 are one leaf of jottings in pencil and black ink, laid out in
%   columns rather than lines. Page 2 gives four equivalent conditions for a
%   class to be primitive — the kernel of Pic -> Hom(A,A*), the theorem-of-the-
%   cube identity, Pic^tau, the image of Ext^1(A,G_m) — with the names of the
%   arguments that join them (Serre, Barsotti-Cartier-Serre). Page 4 carries the
%   Neron-Severi group, the chain algebraic / torsion / numerical, and the chain
%   dim A <= dim Pic A <= dim H^1(A,O_A) <= dim A with the theorems under it.
% - Pages 6 to 10 are a continuous run in blue ink setting out, axiom by axiom,
%   an abstract framework for the theory of abelian varieties: an additive
%   category with a self-duality D, a subgroup N(A) of antisymmetric classes,
%   the cones N^+(A) and N^>(A), an r-linear intersection form I, whence the
%   degree nu(alpha), the coefficients sigma_i, the trace, and the Rosati
%   involution. Its own order is not the archivists'. Page 9 opens on a relative
%   clause continuing the boxed map that closes page 10; page 7 carries "apres
%   d)" in his hand up the left edge, and page 8 is the sheet lettered (d).
%   Against that, the paragraph opening page 8 continues the sigma_i discussion
%   that closes page 7. Both facts are recorded where they occur; neither is
%   resolved here, and the pages are transcribed in the archivists' order.
%
% The hand is fast throughout, and as usual in this fonds the statements carry
% where the connective prose does not (folder 29 batches 8 and 9). Formulas are
% transcribed in full; sentences between them are left \ill{} wherever the leaf
% will not support a reading — including most of the ten-line note that runs
% diagonally across the foot of page 8, which the film does not resolve.
%
% Pass: Opus 5 (claude-opus-5), 2026-09-02 — first pass, unchecked against the pages by a human.
\documentclass[11pt,a4paper]{article}
% Le préambule commun à toutes les transcriptions du fonds Grothendieck.
%
% Il tient en une idée : **l'incertitude se compose, elle ne se résout pas.**
% Une transcription qui lisse un mot douteux a détruit la seule chose pour
% laquelle on la faisait — un lecteur doit pouvoir savoir, à chaque endroit,
% ce qui était sur le feuillet et ce qui a été ajouté. Les six macros qui
% suivent sont donc l'appareil critique, et non de la décoration.
%
% Les mêmes commandes sont reconnues par `scripts/render.mjs`, qui produit la
% vue de lecture du site. Une macro ajoutée ici doit l'être aussi là-bas,
% faute de quoi le rendu échouera bruyamment — ce qui est voulu : un rendu
% silencieusement faux est pire qu'un rendu absent.

\ProvidesPackage{grothendieck}[2026/08/08 Transcriptions du fonds Grothendieck]

\usepackage{amsmath,amssymb,amsthm}
\usepackage{mathrsfs}
\usepackage{stmaryrd}
% Les diagrammes commutatifs sont partout dans ce fonds : c'est la langue
% même de la géométrie algébrique de Grothendieck, et une transcription qui
% les rendrait en prose ne transcrirait rien.
\usepackage{tikz-cd}
% Français par défaut — c'est la langue des feuillets — mais l'anglais est
% chargé aussi : la traduction et le résumé sont des documents anglais, et la
% typographie française (espaces avant les deux-points) y serait une faute.
% Ces documents basculent par \selectlanguage{english} en tête de corps.
\usepackage[english,french]{babel}
\usepackage{fontspec}
\usepackage{geometry}
\geometry{a4paper,margin=25mm,marginparwidth=28mm}
\usepackage{marginnote}
\usepackage{xcolor}
% ulem, not soul. `soul`'s \st and \ul are built on a character-by-character
% scan that XeLaTeX's font machinery breaks: the document fails with an
% undefined control sequence rather than degrading. ulem does the same two jobs
% and is Unicode-engine safe. `normalem` stops it hijacking \emph, which the
% transcriptions use for Grothendieck's own emphasis.
\usepackage[normalem]{ulem}
\usepackage{hyperref}
\hypersetup{colorlinks=true,linkcolor=black,urlcolor=black,pdfborder={0 0 0}}

% --- Métadonnées du lot -----------------------------------------------------
% Reprises telles quelles par le script de rendu, qui les affiche en tête de
% la vue de lecture. Elles fixent ce que le fichier prétend couvrir : sans
% elles, une transcription flotte sans point d'attache dans un fonds de seize
% mille feuillets.

\newcommand{\folder}[1]{\gdef\grothFolder{#1}}
\newcommand{\batch}[1]{\gdef\grothBatch{#1}}
\newcommand{\leaves}[2]{\gdef\grothFirst{#1}\gdef\grothLast{#2}}
\newcommand{\foldertitle}[1]{\gdef\grothFoldertitle{#1}}
\gdef\grothFolder{?}\gdef\grothBatch{?}\gdef\grothFirst{?}\gdef\grothLast{?}\gdef\grothFoldertitle{}

% --- Le résumé --------------------------------------------------------------
%
% Il ouvre la lecture modernisée, sous le titre « Résumé », et il est composé
% en petit. Ce n'est pas un ornement : c'est le seul passage du document qui
% ne soit pas des mathématiques, et un lecteur doit pouvoir le reconnaître
% avant de l'avoir lu — puis le sauter s'il connaît déjà le sujet, ou s'y
% arrêter s'il ne le connaît pas. Le filet qui le suit marque la reprise du
% corps.

\newenvironment{resume}
  {\small\setlength{\parskip}{0.45em}}
  {\par\medskip\hrule\medskip}

% --- Mots-clés ---------------------------------------------------------------
%
% En anglais, en clôture du résumé de la lecture modernisée : le vocabulaire
% moderne sous lequel chercher ce que les feuillets construisent. Le manifeste
% du site (`npm run manifest`) les extrait pour étiqueter la cote entière —
% cette ligne est la source unique des tags, il n'y a pas de fichier à côté.
\newcommand{\keywords}[1]{\par\smallskip\noindent
  {\footnotesize\textsc{Keywords} --- \textit{#1}}\par}

% --- L'appareil critique ----------------------------------------------------

% Le numéro de feuillet, en marge. C'est l'ancre qui permet de régler un
% désaccord sur une formule en regardant *un* feuillet plutôt qu'une chemise
% de six cents. Le site s'en sert aussi pour tourner les pages du fac-similé
% quand on fait défiler la transcription.
\newcommand{\leaf}[1]{%
  \marginnote{\footnotesize\color{gray!70}\textsf{#1}}%
  \label{leaf:#1}\ignorespaces}

% Illisible : on ne devine pas. Un mot inventé qui se lit comme les autres est
% le pire résultat possible.
\newcommand{\ill}{\textcolor{red!60!black}{[\,\ldots\,]}}

% Lecture douteuse : le mot est donné, et signalé comme incertain.
\newcommand{\uncertain}[1]{\uline{#1}}

% Ajout de l'éditeur — une lettre manquante, un mot sous-entendu.
\newcommand{\add}[1]{\textcolor{blue!50!black}{[#1]}}

% Biffé par Grothendieck lui-même. Ce qu'il a rayé fait partie du document :
% une rature dit souvent où la pensée a changé de direction.
\newcommand{\struck}[1]{\sout{#1}}

% Note du transcripteur, sur l'état matériel du feuillet.
\newcommand{\note}[1]{\par\smallskip{\small\color{gray!60!black}\textit{[#1]}}\par\smallskip}

% Note marginale de Grothendieck, distincte du corps du texte.
\newcommand{\marginal}[1]{\par\smallskip\noindent\hspace{1em}%
  {\small\color{gray!60!black}#1}\par\smallskip}

% --- Titre ------------------------------------------------------------------

\AtBeginDocument{%
  \begin{center}
    {\large\bfseries Fonds Alexandre Grothendieck --- cote n\textsuperscript{o}~\grothFolder}\\[2pt]
    {\itshape\grothFoldertitle}\\[4pt]
    {\small feuillets \grothFirst--\grothLast\ (lot \grothBatch)}\\[2pt]
    {\footnotesize\color{gray!60!black}Transcription établie d'après le fac-similé numérisé
     par l'Université de Montpellier.\\
     Les lectures incertaines sont soulignées, les illisibles notées [\,\ldots\,],
     les ajouts entre crochets.}
  \end{center}
  \vspace{1em}}

\endinput


\folder{50}
\batch{1}
\pages{1}{10}
\dating{[à partir de 1965]}
\watermark{Édition de démonstration}
\foldertitle{Sur bouquin Lang : notes manuscrites (s.d.)}

\begin{document}

\section*{Sur bouquin Lang}

\note{titre de sa main, au crayon, dans l'angle supérieur droit du premier
feuillet — une chemise grise qui ne porte rien d'autre. Le feuillet ne reçoit
donc pas de \texttt{page}, conformément à l'usage suivi depuis le dossier 29 :
le saut dans la numérotation est la trace qu'il a été passé. « bouquin Lang »
est son mot pour le livre de Lang, \emph{Abelian Varieties} ; en tête de la
page 10 il écrit « livre de Lang / VA »}

\page{2}

\note{feuillet de travail sans ordre linéaire : les blocs sont posés en
colonnes, et la transcription les suit de haut en bas. La moitié supérieure
fixe les objets, la moitié inférieure énonce quatre conditions équivalentes et
nomme les arguments qui les joignent}

\note{en tête, un départ abandonné — « Nérm( » — relié par un long trait à
« $A \longrightarrow$ » écrit au-dessus}

\[
  \underline{\underline{\text{Nérm}}}_{A/S} \longrightarrow
  \underline{\underline{\mathrm{Hom}}}_{S\text{-}\mathrm{gr}}(A, A^{*})
\]

\note{« Nérm » est la lecture littérale des quatre lettres ; il s'agit du
foncteur de Néron--Severi, comme l'impose la flèche vers
$\underline{\underline{\mathrm{Hom}}}_{S\text{-}\mathrm{gr}}(A,A^{*})$ — au
dossier 49 il l'abrège « Nér.-Sév. ». Le double soulignement, ici et partout
sur ce feuillet, marque le foncteur}

\begin{tikzcd}
A \times A \arrow[d, "\pi"] \\
A
\end{tikzcd}

\note{deux signes minuscules surmontent les deux facteurs de $A \times A$ et ne
se laissent pas lire ; à droite du diagramme, $x \longmapsto x - s$}

À gauche du diagramme : $(T_{s})_{*}(\xi) - \xi$ ; au-dessous,
\struck{$A \longrightarrow \uncertain{D}^{*}$}, puis $V$, $s \times A$,
$s_{*}\xi - \xi$.

\[
  \boxed{\ \pi^{*}(\xi) - \mathrm{pr}_{1}(\xi) - \mathrm{pr}_{2}(\xi)\ }
\]

\note{les deux projections sont écrites sans étoile, à la différence de
$\pi^{*}$ ; il en va de même à la ligne (i$'$) ci-dessous}

(i)\quad $\xi$ dans le noyau \ill{} $s^{*}\xi$ \ill{} de
$\underline{\underline{\mathrm{Pic}}}_{A/S} \longrightarrow
\underline{\underline{\mathrm{Hom}}}(A, A^{*})$, i.e. $s\xi = \xi$

(i$'$)\quad $\pi^{*}(\xi) - \mathrm{pr}_{1}(\xi) - \mathrm{pr}_{2}(\xi) = 0$

(ii)\quad $\xi \in \underline{\underline{\mathrm{Pic}}}^{\tau}
\struck{\ill{}}(A/S)$

(iii)\quad $\xi \in \mathrm{Im}\ \mathrm{Ext}^{1}_{\uncertain{S}}(A,
\mathbb{G}_{m})$

\note{l'indice de $\mathrm{Ext}^{1}$ est tracé comme son $\xi$ ; c'est la base
$S$ que demande le sens}

En marge de ces quatre lignes : $s.D \sim D$.

(i)\quad \uncertain{$\Longrightarrow$}\quad (i$'$)\quad banal

(i$'$)\quad $\Longleftrightarrow$\quad (iii)\quad par le raisonnement
\uncertain{giraudique} (de Serre

(ii)\quad $\Longrightarrow$\quad (i)\quad par la propriété additive de
$\underline{\underline{\mathrm{Pic}}}^{\tau}$

(i)\quad $\Longrightarrow$\quad (ii)\quad c'est le th. de
Barsotti--Cartier--Serre

\note{le signe de la première ligne est petit et mis entre parenthèses sur la
feuille ; il porte une pointe à droite, mais une équivalence n'est pas exclue.
À la deuxième ligne, « giraudique » est écrit en interligne et relié par un
trait à la parenthèse ouvrante de « (de Serre », qu'il ne referme pas}

\page{4}

\note{même feuillet de travail que la page 2, déchiré en haut ; les blocs y sont
encore disposés en colonnes}

En tête : $\mathrm{N.S.}$

additivité de $\underline{\underline{\mathrm{Pic}}}^{\tau}$, on \ill{}
spécial : $dA = A$

\note{une première esquisse de la chaîne — « torsion », un « $\equiv$ »,
« équivalence numérique » — occupe la colonne de gauche ; elle est reprise plus
bas sous forme encadrée, et c'est la même chaîne écrite deux fois}

[critère de Weil.--]\quad \ill{}\quad \ill{}\quad variations\quad
\ill{} pour les VA (bouquin Lang)

\marginal{dû à M. \uncertain{Matsusaka}}

\note{trois lignes de commentaire accompagnent une flèche montant vers
« N.S. » ; seuls « [critère de Weil.-- », « variations » et « pour les VA
(bouquin Lang) » se laissent lire. La mention entourée « dû à M. \ill{} » est
une lecture douteuse : le nom, en neuf ou dix lettres, finit en « -aka »}

\begin{tikzcd}
\text{algébrique} \arrow[d, Rightarrow] \\
\text{torsion} \arrow[d, Rightarrow] \\
\text{numérique}
\end{tikzcd}

\note{encadré ; entre « torsion » et « numérique » il pose un « $\equiv$ » et
un trait}

\[
  \dim A \;\leq\; \dim \mathrm{Pic}\,A \;\leq\;
  \dim H^{1}(A, \mathcal{O}_{A}) \;\leq\; \dim A
\]

\note{entouré d'un trait ; trois accolades le commentent — « théorie des
diviseurs non dégénérés » sous la première inégalité, « déf. de Picard » sous
la deuxième, « Hopf-Borel » sous la troisième. Au-dessus, une ligne en partie
biffée : « \struck{\ill{}} définition de Picard ; \ill{} »}

$A^{**}$\quad $A$\,.\quad $A^{*}$

\note{la ligne se poursuit par deux symboles plus petits qui ne se laissent pas
lire}

\page{6}

\textbf{Th.}\quad Soit $\xi \in N^{>}(A)$, considérons l'involution

\[
  \alpha \rightsquigarrow \alpha' = \struck{\ill{}}\ \varphi_{\xi}^{-1}\,
  {}^{t}\alpha\,\varphi_{\xi}
\]

de $\mathrm{Hom}(A,A)_{\mathbb{Q}}$, alors \ill{}

\[
  \mathrm{tr}(\alpha\beta') = \frac{r}{I(\xi, \ldots, \xi)}\,
  I(\xi, \ldots, \xi, D_{\xi}(\alpha\beta'))
\]

d'où

\[
  \mathrm{tr}\,\alpha\alpha' = \frac{r}{I(\xi, \ldots, \xi)}\,
  I(\xi, \ldots, \xi, 2\alpha^{*}(\xi))
\]

\textbf{Corollaire}

\note{le corollaire n'est pas énoncé : le feuillet s'arrête sur le mot,
souligné, et tout le reste de la page est blanc}

\page{7}

\marginal{après d)}

\note{écrit dans l'angle supérieur gauche, la feuille tournée d'un quart de
tour : c'est son propre repère d'ordre, et il place ce feuillet après celui qui
porte (d), c'est-à-dire après la page 8}

(e)\quad Pour tout $A \in \mathrm{Ob}\,\mathcal{A}$, \ill{} une forme
$r$-fois \emph{linéaire}

\[
  (\xi_{1}, \ldots, \xi_{r}) \rightsquigarrow I(\xi_{1}, \ldots, \xi_{r})
  \in \struck{\mathbb{Q}}\ \mathbb{Z}
\]

\note{en interligne, relié au corps du texte par un trait : « \uncertain{un
entier} $r = \dim(A)$ »}

2)\quad On suppose $I$ \struck{\ill{}} \uncertain{str.\ pos.} sur $N^{>}(A)$,
i.e. $\xi_{1}, \struck{\ill{}}\ \xi_{r} \in N^{+}(A) \Longrightarrow
I(\xi_{1}, \ldots, \xi_{r}) > 0$.

\note{les deux abréviations se lisent « sh. po. » sur la feuille ; c'est le
sens qui impose « strictement positive »}

3)\quad On \ill{} pour \ill{} $u : A \longrightarrow B$ \ill{} avec
$\dim(A) = \dim(B)$, \struck{\ill{}} il existe un entier \struck{\ill{}}
$\nu(u)$ \ill{}

\[
  I(u^{*}(\xi_{1}), \ldots, u^{*}(\xi_{r})) = \nu(u)\,
  I(\xi_{1}, \ldots, \xi_{r})
\]

\ill{} $\xi_{1}, \ldots, \xi_{r}$ [\uncertain{ce qui le détermine} \ill{}
\uncertain{$\mathbb{Q}$-isomorphisme} \ill{}]. \ill{} en particulier \ill{}
$\xi \in N(B)$ \ill{} et $I(\xi, \ldots, \xi) \neq 0$ \ill{} pour un certain
$\xi$ \ill{}

\[
  \nu(\alpha) = \frac{1}{I(\xi, \ldots, \xi)}\,
  I(\alpha^{*}(\xi), \ldots, \alpha^{*}(\xi))
\]

$=$ \emph{forme de degré} $2r$ \emph{en} $\alpha$, à coefficients
\uncertain{entiers} [\ill{} si $A = B$], donc

\[
  \nu(\alpha + n\cdot 1_{A}) = n^{2r} + \sigma_{1}(\alpha)\,n^{2r-1} +
  \cdots + \sigma_{2r}(\alpha)
\]

et on posera $\sigma_{1}(\alpha) = \mathrm{tr}\,\alpha$, tandis qu'on a
$\sigma_{2r}(\alpha) = \nu(\alpha)$, noté aussi $\det \alpha$.

\page{8}

de façon générale \struck{$\varphi_{\xi}(\alpha)$} les $\sigma_{i}(\alpha)$
sont des formes de degré $2i$ en $\alpha$, \uncertain{etc.} en particulier,
\ill{}

\[
  \mathrm{tr}(\alpha) = \frac{r}{I(\xi, \ldots, \xi)}\,
  I(\xi, \ldots, \xi, D(\alpha))
\]

\note{ce paragraphe prolonge les dernières lignes de la page 7, où les
$\sigma_{i}$ viennent d'être introduits, alors que le même feuillet porte plus
bas l'axiome (d), qui précède (e) : voir la note de la page 7. Les deux faits
sont consignés là où ils se trouvent et la question n'est pas tranchée ici}

(d)\quad Pour tout $A \in \mathrm{Ob}\,\mathcal{A}$, \ill{} dans $N(A)$ des
parties

\[
  N^{+}(A),\quad N^{>}(A) \;\subset\; N(A)
\]

dépendant fonctoriellement de $A$, et pour les \struck{isomorphismes},
\uncertain{vérifiant} les propriétés suivantes :

\note{l'exposant du premier $N$ est repassé deux fois et se lit tour à tour
$\tau$, $+$ et $\neq$ ; les lignes suivantes distinguent nettement $N^{+}$ et
$N^{>}$, et la transcription suit cette distinction}

a)\quad $N^{+}(A) + N^{+}(A) \subset N^{+}(A)$, \struck{\ill{}} $n > 0
\Longrightarrow n\,N^{+}(A) \ill{}$ ; kif kif pour $N^{>}(A)$

b)\quad $N^{>}(A)$ engendre $N(A)$

\struck{c) $\xi \in N^{+}(A) \Longrightarrow I(\xi \ill{}) > 0$ \ill{}
$N^{+}(A)$}

\struck{i.e. $\xi_{1}, \ldots, \xi_{r} \in N^{>}(A) \Longrightarrow
I(\xi_{1}, \ldots, \xi_{r}) > 0$}

c)\quad $\xi \in N^{>}(A) \Longrightarrow \varphi_{\xi}$ soit un
$\mathbb{Q}$-isom.

\marginal{\ill{}}

\note{une note d'une dizaine de lignes court en diagonale au bas du feuillet,
de la même encre et entre deux longs traits ; on y reconnaît $N^{+}(A)$,
$n > 0$, « diviseurs », « l'image \ill{} » et $N^{>}(B)$, sans que la suite se
laisse lire à la définition que porte le film}

\page{9}

qui d'ailleurs s'annule sur $\mathrm{Im}\,N(A) + \mathrm{Im}\,N(B)$

\textbf{Propriétés fondamentales}\quad 1) On \uncertain{trouve} ainsi un
\emph{isom.}

\[
  N(A\times B) \big/ N(A)\times N(B) \;\xrightarrow{\ \sim\ }\;
  \mathrm{Hom}(A, D(B))
\]

\marginal{purement \emph{formel}}

\note{cette note marginale est au crayon, d'une autre pointe que le corps du
feuillet}

2)\quad Pour tout $A$, il existe des $\xi \in N(A)$ telles que
$\varphi_{\xi} : A \xrightarrow{\ \sim\ } D(A) \otimes \mathbb{Q}$,
\struck{\ill{}} i.e. $\varphi_{\xi} : A \longrightarrow D(A)$ soit un
$\mathbb{Q}$-isom.,

\note{trois traits diagonaux barrent un « N.B. » entre crochets qui invoque
« \ill{} les th\ill{} des var. ab. pour $\mathbb{Q}$ » et ne se laisse pas lire
davantage}

Cela permet donc de définir une involution sur
$\mathrm{Hom}_{\mathbb{Q}}(A,A)$ par

\[
  s_{\xi}(\alpha) = \varphi_{\xi}^{-1}\,{}^{t}\alpha\,\varphi_{\xi}
\]

C'est là une \emph{involution} \struck{\ill{}} de \emph{l'anneau}
$\mathrm{Hom}(A,A)$

Posons pour $\alpha, \beta \in \mathrm{Hom}_{\mathbb{Q}}(A,A)$ :

\[
  \text{(1)}\qquad D_{\xi}(\alpha,\beta) = (\alpha+\beta)^{*}(\xi) -
  \alpha^{*}(\xi) - \beta^{*}(\xi) \in N(A)
\]

i.e.

\[
  \varphi_{D_{\xi}(\alpha,\beta)} = {}^{t}(\alpha+\beta)\,\varphi_{\xi}\,
  (\alpha+\beta) - {}^{t}\alpha\,\varphi_{\xi}\,\alpha -
  {}^{t}\beta\,\varphi_{\xi}\,\beta
\]

i.e.

\[
  \text{(1$'$)}\qquad \varphi_{D_{\xi}(\alpha,\beta)} =
  {}^{t}\beta\,\varphi_{\xi}\,\alpha + {}^{t}\alpha\,\varphi_{\xi}\,\beta
  = \varphi_{\xi}(\alpha'\beta + \beta'\alpha)
\]

\note{dans $(\alpha+\beta)$ le signe est tracé comme un petit anneau et non
comme une croix, ici et à la ligne suivante ; c'est l'identité (1$'$) qui
impose la somme}

(où on \uncertain{pose} $\alpha' = s_{\xi}(\alpha)$) ; \ill{} posons aussi

\[
  \text{(2)}\qquad D_{\xi}(\alpha) = D_{\xi}(\alpha, \mathrm{id}_{A}) =
  \varphi_{\xi}(\alpha' + \alpha)
\]

d'où

\[
  \text{(3)}\qquad D_{\xi}(\alpha,\beta) = D_{\xi}(\alpha'\beta) =
  D_{\xi}(\beta'\alpha), \qquad D_{\xi}(\alpha) = D_{\xi}(\alpha')
\]

\page{10}

\marginal{livre de Lang / VA}

\note{au crayon, en tête du feuillet, d'une autre pointe que le corps du texte}

(a)\quad \emph{Catégorie additive} $\mathcal{A}$

(b)\quad \emph{Autodualité} : $D : \mathcal{A}^{\circ} \longrightarrow
\mathcal{A}$ (munie d'un isomorphisme $\kappa : \mathrm{id}_{\mathcal{A}}
\longrightarrow D \circ D$ satisfaisant la condition de compatibilité
habituelle).

Soit donc

\[
  B(A,B) = \mathrm{Hom}(A, D(B)) \;\xrightarrow[\ \simeq\ ]{s}\;
  \mathrm{Hom}(B, D(A))
\]

c'est un bifoncteur contravariant en $A$, $B$. Le \ill{}

\note{trois mots suivent « Le » à la ligne suivante et ne se laissent pas lire}

(c)\quad Pour tout $A \in \mathrm{Ob}\,\mathcal{A}$, un sous-groupe

\[
  \boxed{\ N(A) \subset B(A,A)\ }
\]

\note{« antisymétrique » est écrit au-dessus de l'encadré}

dépendant de $A$ de façon \struck{\ill{}} contravariante.

[N.B. $N$ n'est pas additif en $A$ \ill{}]

Si $\xi \in N(A)$, on \ill{} $\varphi_{\xi} : A \longrightarrow D(A)$ le
morphisme qu'il définit.

Soit $\xi \in N(A\times B)$, il définit

\[
  \varphi_{\xi} : A\times B \longrightarrow D(A)\times D(B)
\]

qui est antisymétrique, donc \emph{nécessairement} de la forme

\[
  \begin{pmatrix}
    \alpha_{\xi} & {}^{t}\lambda_{\xi} \\
    -\lambda_{\xi} & \beta_{\xi}
  \end{pmatrix},
  \qquad \lambda_{\xi} \in \mathrm{Hom}(A, D(B)) = B(A,B),\quad
  \alpha_{\xi} \in B(A,A),\quad \beta_{\xi} \in B(B,B)
\]

on obtient ainsi une \ill{}

\[
  \boxed{\ \xi \longmapsto \lambda_{\xi} : N(A\times B) \longrightarrow
  \mathrm{Hom}(A, D(B)) = B(A,B)\ }
\]

\note{la page 9 s'ouvre sur la relative qui prolonge cet encadré — « qui
d'ailleurs s'annule sur $\mathrm{Im}\,N(A) + \mathrm{Im}\,N(B)$ »}

\end{document}
